The linear form $\tau_{-1}$  provides a substitute on
abstract pseudodif\/ferential operators with order outside the discrete set $-{\rm Sd}$, for the canonical trace TR
def\/ined on classical pseudodif\/ferential
operators, whose order lies outside  the set $[{-}n,+\infty[\cap \Z$.
\begin{proposition}\label{proposition7}
The linear\footnote{Meaning by this that it acts on  linear
  combinations of operators whose order does not lie in $-{\rm Sd}$.}  map
$\tau_{-1}$  vanishes on  commutators, whose order does not  lie in~$-{\rm Sd}$:
\[
\left({\rm ord}(A)+{\rm ord}(B)\notin -{\rm Sd}\right) \ \Longrightarrow \
\tau_{-1}\left([A,B]\right)=0.
\]
\end{proposition}

\begin{proof}
By (\ref{eq:taujDbracket}) we have
\begin{gather*}
\tau_{-1}([A,B]) = \sum_{n> 0} (-1)^{n+1}\frac{\tau_{n-1}^{\vert D\vert} \left(
    A L^n(B)\right)}{n!}
 =  \sum_{m\geq  0} (-1)^{m}\frac{\tau_{m}^{\vert D\vert} \left(
    A L^{m+1}(B)\right)}{(m+1)!}.
\end{gather*}
The orders of the  operators
  $  A  L^{m+1}(B)$, which lie  in ${\rm ord}(A)+ {\rm ord}(B)-\Z_{\geq 0}$  do not lie
  in $-{\rm Sd}$ since   ${\rm ord}(A)+ {\rm ord}(B)$ does not and since by assumption (T) we have the inclusion: $ -{\rm Sd}\subset
-{\rm Sd} -\Z_{\geq 0}$. Hence,
 the terms $\tau_{m}^{\vert D\vert} \left(
    A\,L^{m+1}(B)\right)$, $m\in \Z_{\geq 0}$ in the sum vanish and
$\tau_{-1}([A,B])=0$. \end{proof}

 Thus, the linear map $\tau_{-1}$ def\/ines a ``trace''  on the set of
operators, whose order  lies outside the set $-{\rm Sd}$  in the same way that
Kontsevich and Vishik's canonical trace TR
def\/ines a~``trace'' on non-integer order classical pseudodif\/ferential
operators.
